
%% --------------------------------------------------------
\subsection*{Поляризація випромінювання}
%% --------------------------------------------------------


При розгляді поляризації треба зважати на векторну природу польових функцій, зокрема, напруженості електричного поля. Зауважимо, що для фіксованого
напрямку випромінювання досить розглядати лише електричне поле; воно визначає також і вектор напруженості магнітного поля.

Розглянемо поле випромінювання, що поширюється в напрямку осі \(x_3\), з напруженістю електричного поля \(\Efield = (E_1, E_2, 0)\), причому будемо
вважати, що поле є суперпозицією монохроматичних хвиль \(\sim e^{-i\omega t}\) з близькими частотами. Напрямок вектора \(\Efield\), що є ортогональним
напрямку поширення хвилі, визначає поляризацію (рис.~\ref{pic:E-direction}).

%---------------------------------------------------------
\begin{SCfigure}[][h!]\centering
	\begin{tikzpicture}[scale=1.5]
    	      \draw[-{Stealth[length=2.8pt, bend]},thin,black!55]
          ({-1.55*1},0)--({1.55*1},0)
          node[right,font=\tiny,inner sep=1pt]{$X_1$};
      \draw[-{Stealth[length=2.8pt, bend]},thin,black!55]
          (0,{-1.55*1})--(0,{1.55*1})
          node[above,font=\tiny,inner sep=1pt]{$X_2$};
            \def\a{45}
            \def\l{1.45}
        \draw[-{Latex[length=6pt]}, red, thick] (0,0) -- ++(\a:\l) coordinate (E) node[pos=1, anchor=south] {$\Efield$};
        \draw[densely dashed] (E) -- node[anchor={180}] {$E_2$} ++(0,{-\l*sin(\a)});
        \draw[densely dashed] (E) --  node[anchor={south}] {$E_1$}  ++({-\l*cos(\a)}, 0);
	\end{tikzpicture}
\caption{Напрямок вектора \(\Efield\) визначає поляризацію хвилі.}
\label{pic:E-direction}
\end{SCfigure}
%---------------------------------------------------------


Для випадкового поля \(\Efield\) введемо тензор поляризації:

\begin{equation}
\rho_{ij} = \frac{\langle E_i E_j^{*} \rangle}{\langle |E_1|^2 \rangle + \langle |E_2|^2 \rangle}, \quad i, j = 1, 2
\label{eq:polarization_tensor}
\end{equation}

Очевидно, матриця \( \|\rho\|\) є ермітовою, а її шпур (слід) дорівнює одиниці:

\begin{equation}\label{eq:hermitian}
\rho_{ji} = \rho_{ij}^{*}, \quad \rho_{11} + \rho_{22} = 1
\end{equation}

Для цілком поляризованого випромінювання маємо \(E_i = N_i f(t)\), де \(f\) може бути випадковою функцією, але вектор \(N_i\), що визначає напрямок
поля, є фіксованим. В цьому випадку \(\rho_{ij}\) пропорційне \(N_i N_j^{*}\) і:

\begin{equation}
\det \|\rho\| = 0
\end{equation}

У випадку природного світла присутні усі можливі поляризації, що є рівноправними, причому \(\langle E_1 E_2^{*} \rangle = 0\), \(\langle |E_1|^2 \rangle =
\langle |E_2|^2 \rangle\). В цьому разі:

\begin{equation*}
\rho_{ij} = \frac{1}{2} \delta_{ij}
\end{equation*}

В загальному випадку визначник матриці \(\|\rho\|\) є:

\begin{equation*}
\det \|\rho\| = \rho_{11} \rho_{22} - |\rho_{12}|^2 = \frac{\langle |E_1|^2 \rangle \langle |E_2|^2 \rangle - |\langle E_1 E_2^{*} \rangle|^2}{
\langle |E_1|^2 \rangle + \langle |E_2|^2 \rangle}
\end{equation*}

Оскільки \(|\langle E_1 E_2^{*} \rangle|^2 \leqslant  \langle |E_1|^2 \rangle \langle |E_2|^2 \rangle\), цей визначник додатний:

\begin{equation*}
\det \|\rho\| \geqslant 0
\end{equation*}

З іншого боку, очевидно:

\begin{equation*}
\det \|\rho\| \leqslant
\frac{\langle |E_1|^2 \rangle \, \langle |E_2|^2 \rangle}
{\left(\langle |E_1|^2 \rangle + \langle |E_2|^2 \rangle\right)^2}
\leqslant \frac{1}{4}
\end{equation*}

Оскільки область значень \(\det \|\rho\| \) є відрізком \([0, 1/4]\), можна записати:

\begin{equation}
\det \|\rho\| = \frac{1}{4} \left( 1 - p^2 \right),\quad  \text{де}\  0 \leqslant  p \leqslant 1
\label{eq:polarization_degree}
\end{equation}
%
де параметр \(p\) називають \emph{ступінню поляризації}\index{Ступінь поляризації}; він приймає значення від 0 (неполяризоване світло) до 1.

Завдяки властивості~\eqref{eq:hermitian} (ермітовості), власні числа \(\lambda_1, \lambda_2\) матриці \(\rho_{ij}\) є дійсними, причому \(\det \|\rho\| = \lambda_1 \lambda_2\),
\(\lambda_1 + \lambda_2 = 1\). Зважаючи на~\eqref{eq:polarization_degree}, добуток власних чисел додатний, тобто вони мають один знак. Звідси легко
бачити, що ці числа є додатними. Якщо матриця $\|\rho \|$ --- дійсна, її можна привести до діагонального виду  $\| \rho \| = \mathrm{diag}(\lambda_{\min}, \lambda_{\max})$ за допомогою просторових поворотів. Частку \(\lambda_{\min}/\lambda_{\max}\) називають \emph{коефіцієнтом деполяризації}\index{Коефіцієнт деполяризації}.

%% --------------------------------------------------------
\subsection*{Параметри Стокса}
%% --------------------------------------------------------

Виходячи з властивостей тензора поляризації, поданих формулою~\eqref{eq:polarization_tensor}, його компоненти \(\rho_{ij}\) можна виразити через три
незалежні дійсні \emph{параметри Стокса}\index{Параметри Стокса} \(\xi_1, \xi_2, \xi_3\), які приймають значення від $-1$ до $1$:

\begin{equation*}
\rho_{11} = \frac{1}{2} \left( 1 + \xi_3 \right), \quad \rho_{22} = \frac{1}{2} \left( 1 - \xi_3 \right), \quad \rho_{12} = \frac{1}{2} \left( \xi_1 -
i\xi_2 \right) = \rho_{21}^{*}
\end{equation*}

Очевидно:

\begin{equation*}
\det \|\rho\| = \frac{1}{4} \left( 1 - \xi_1^2 - \xi_2^2 - \xi_3^2 \right)
\end{equation*}

\begin{equation*}
p^2 = \xi_1^2 + \xi_2^2 + \xi_3^2, \quad 0 \leqslant \xi_1^2 + \xi_2^2 + \xi_3^2 \leqslant 1
\end{equation*}

Розглянемо стани поляризації для деяких полів, що відповідають різним параметрам Стокса.

\begin{itemize}

\item Лінійна поляризація під кутом $45^{\circ}$, тобто \(E_1 = \pm E_2\):

\begin{equation*}
\rho_{11} = \rho_{22} = \frac{1}{2} \rightarrow \xi_3 = 0; \quad \text{Im} \, \rho_{12} = 0 \rightarrow \xi_2 = 0;
\end{equation*}

Таким чином:

\begin{equation*}
\xi_1 = \pm 1, \quad \xi_3 = \xi_2 = 0
\end{equation*}

%---------------------------------------------------------
\begin{SCfigure}[][h!]\centering
	\begin{tikzpicture}[scale=1.5]
    	      \draw[-{Stealth[length=2.8pt, bend]},thin,black!55]
          ({-1.55*1},0)--({1.55*1},0)
          node[right,font=\tiny,inner sep=1pt]{$X_1$};
      \draw[-{Stealth[length=2.8pt, bend]},thin,black!55]
          (0,{-1.55*1})--(0,{1.55*1})
          node[above,font=\tiny,inner sep=1pt]{$X_2$};
            \def\a{45}
            \def\l{1.45}
        \draw[-{Latex[length=6pt]}, red, thick] (0,0) -- ++(\a:\l) ;
        \draw[-{Latex[length=6pt]},dashed,  red, thick] (0,0) -- ++(180+\a:\l) ;
        \draw[densely dashed] (E) -- node[anchor={180}] {$E_2$} ++(0,{-\l*sin(\a)});
        \draw[densely dashed] (E) --  node[anchor={south}] {$E_1$}  ++({-\l*cos(\a)}, 0);
	\end{tikzpicture}
\caption{Лінійна поляризація під кутом $45^{\circ}$ (випадок $E_1 = E_2$).}
\label{pic:E-lin45}
\end{SCfigure}
%---------------------------------------------------------

\item  Колова поляризація \(E_1 = A e^{-i\omega t}, \quad E_2 = \pm i E_1\).

Як видно з поведінки дійсних частин компонент, вектор \(\Efield\) обертається в площині \(X_1 -
X_2\). Маємо:

\begin{equation*}
\rho_{11} = \frac{1}{2}, \quad \rho_{22} = \frac{1}{2}, \quad \rho_{12} = \mp i
\end{equation*}

Звідси:

\begin{equation*}
\xi_1 = \xi_3 = 0, \quad \xi_2 = \pm 1
\end{equation*}

%---------------------------------------------------------
\begin{SCfigure}[1.5][h!]\centering
	\begin{tikzpicture}[scale=1.5]
    	      \draw[-{Stealth[length=2.8pt, bend]},thin,black!55]
          ({-1.55*1},0)--({1.55*1},0)
          node[right,font=\tiny,inner sep=1pt]{$X_1$};
            \def\a{45}
            \def\l{1.2}
        \draw[-{Latex[length=6pt]}, red, thick] (0,0) -- ++(\a:\l) coordinate (E) node[pos=1, anchor=south] {$\Efield$};
      \draw[-{Stealth[length=2.8pt, bend]},black!55]
          (0,{-1.55*1})--(0,{1.55*1})
          node[above,font=\tiny,inner sep=1pt]{$X_2$};
          \draw[-{Stealth[length=4.5pt, bend]},thick, red]
              (\l,0)arc[start angle=0,end angle=360,
                         radius=\l];
	\end{tikzpicture}
\caption{Колова поляризація (випадок $E_2 = iE_1$, обертання вектора $\Efield$ проти годинникової стрілки).}
\label{pic:E-circle}
\end{SCfigure}
%---------------------------------------------------------

\item Лінійна поляризація вздовж однієї з осей \(E_1 \neq 0, E_2 = 0\) або \(E_2 \neq 0, E_1 = 0\):

\begin{equation*}
\rho_{11} = 1, \quad \rho_{22} = 0 \rightarrow \xi_3 = 1,
\end{equation*}

або

\begin{equation*}
\rho_{11} = 0, \quad \rho_{22} = 1 \rightarrow \xi_3 = -1,
\end{equation*}

при цьому \(\rho_{12} = 0 \rightarrow \xi_1 = \xi_2 = 0\).

%---------------------------------------------------------
\begin{SCfigure}[][h!]\centering
	\begin{tikzpicture}[scale=1.5]
    	      \draw[-{Stealth[length=2.8pt, bend]},thin,black!55]
          ({-1.55*1},0)--({1.55*1},0)
          node[right,font=\tiny,inner sep=1pt]{$X_1$};
      \draw[-{Stealth[length=2.8pt, bend]},thin,black!55]
          (0,{-1.55*1})--(0,{1.55*1})
          node[above,font=\tiny,inner sep=1pt]{$X_2$};
            \def\a{0}
            \def\l{1.4}
        \draw[-{Latex[length=6pt]}, red, thick] (0,0) -- ++(\a:\l) ;
        \draw[-{Latex[length=6pt]},dashed,  red, thick] (0,0) -- ++(180+\a:\l) ;
	\end{tikzpicture}
\caption{Лінійна поляризація вздовж $X_1$ (випадок $E_2 = 0$).}
\label{pic:E-lin0}
\end{SCfigure}
%---------------------------------------------------------

Таким чином:

\begin{equation*}
\xi_1 = \xi_2 = 0, \quad \xi_3 = \pm 1
\end{equation*}


\item Еліптична поляризація --- загальний випадок: \(E_2 = A e^{i\delta} E_1\)
з довільною амплітудою \(A\) та зсувом фаз \(\delta\). Тоді кінець вектора
\(\Efield\) описує еліпс у площині \(X_1\!-\!X_2\), а компоненти тензора
поляризації є нетривіальними:

\begin{equation*}
\rho_{11} = \frac{1}{1 + A^2}, \quad \rho_{22} = \frac{A^2}{1 + A^2}, \quad
\rho_{12} = \frac{A e^{-i\delta}}{1 + A^2}
\end{equation*}

Звідси всі три параметри Стокса, взагалі кажучи, відмінні від нуля:

\begin{equation*}
\xi_3 = \frac{1 - A^2}{1 + A^2}, \quad
\xi_1 = \frac{2A\cos\delta}{1 + A^2}, \quad
\xi_2 = \frac{2A\sin\delta}{1 + A^2}
\end{equation*}

причому \(\xi_1^2 + \xi_2^2 + \xi_3^2 = 1\). Лінійна (\(\delta = 0,\, \pi\))
та кругова (\(A = 1,\, \delta = \pm\pi/2\)) поляризації є граничними випадками.
%---------------------------------------------------------
\begin{SCfigure}[][h!]\centering
	\begin{tikzpicture}[scale=1.5]
    	      \draw[-{Stealth[length=2.8pt, bend]},thin,black!55]
          ({-1.55*1},0)--({1.55*1},0)
          node[right,font=\tiny,inner sep=1pt]{$X_1$};
            \def\a{45}
            \def\l{1.2}
            \draw[-{Latex[length=6pt]}, red, thick] (0,0) -- ({\l*cos(\a)},{0.65*\l*sin(\a)}) coordinate (E) node[pos=1, anchor=south] {$\Efield$};
      \draw[-{Stealth[length=2.8pt, bend]},black!55]
          (0,{-1.55*1})--(0,{1.55*1})
          node[above,font=\tiny,inner sep=1pt]{$X_2$};
          \draw[-{Stealth[length=4.5pt, bend]},thick, red]
              (\l,0)arc[start angle=0,end angle=360,
                         x radius=\l, y radius=0.65*\l];
	\end{tikzpicture}
\caption{Еліптична поляризація.}
\label{pic:E-ellipse}
\end{SCfigure}
%---------------------------------------------------------

\end{itemize}

%% --------------------------------------------------------
\subsection*{Сфера Пуанкаре}
%% --------------------------------------------------------

Параметри Стокса \((\xi_1, \xi_2, \xi_3)\) зручно інтерпретувати як декартові координати точки на сфері радіуса \(p = 1\) (ступінь поляризації). Стан поляризації цілком поляризованого випромінювання (\(p = 1\)) однозначно задається точкою на одиничній сфері \(\xi_1^2 + \xi_2^2 + \xi_3^2 = 1\) --- \emph{сфері Пуанкаре}\index{Сфера Пуанкаре} (рис.~\ref{pic:Stoke})\footnote{Сфера Пуанкаре перетворює задачу про поляризацію на \emph{задачу обертань цієї сфери}: дія кожного поляризаційного елемента (фазова пластинка, обертач тощо) описується поворотом сфери на певний кут навколо певної осі. Це дозволяє наочно й без громіздких обчислень передбачати результат дії будь-якої оптичної системи на поляризоване світло.}. Лінійні та кругові поляризації займають на ній особливі положення (полюси і точки екватора), тоді як еліптична поляризація відповідає всім решта точкам поверхні. Частково поляризоване світло зображується точкою \emph{всередині} сфери, а неполяризоване --- її центром.


Вигляд поляризації у площині \(X_1\!-\!X_2\) повністю визначається двома кутами: кутом нахилу великої осі \(\psi \in (-90^{\circ}, 90^{\circ})\) та кутом еліптичності \(\chi \in (-45^{\circ}, 45^{\circ})\), де \(\tg\chi = b/a\) --- відношення малої півосі до великої, зі знаком що визначає напрямок обертання (\(\chi > 0\) --- проти годинникової стрілки). Зв'язок цих кутів з параметрами Стокса:

\begin{equation*}
\xi_1 = \cos 2\chi \cos 2\psi, \quad
\xi_2 = \sin 2\chi, \quad
\xi_3 = \cos 2\chi \sin 2\psi
\end{equation*}

Обернені формули:

\begin{equation*}
\psi = \tfrac{1}{2}\arctan\frac{\xi_1}{\xi_3}, \quad
\chi = \tfrac{1}{2}\arcsin\xi_2
\end{equation*}

Таким чином, кути \(\psi\) і \(\chi\) грають роль довготи і широти на сфері Пуанкаре. Граничні випадки: \(\chi = 0\) --- лінійна поляризація під кутом \(\psi\); \(\chi = \pm 45^{\circ}\) --- кругова.


%---------------------------------------------------------
\begin{figure}[ht!]\centering
% ================================================================
%  ТЕОРІЯ (з тексту):
%    E2 = A*exp(i*delta)*E1
%    xi3 = (1-A^2)/(1+A^2)
%    xi1 = 2A*cos(delta)/(1+A^2)
%    xi2 = 2A*sin(delta)/(1+A^2)
%
%  Еліпс у площині X1-X2:
%    кут нахилу великої осі:  psi = 0.5 * atan2(xi1, xi3)
%    кут еліптичності:        chi = 0.5 * asin(xi2 / p)
%    p = sqrt(xi1^2+xi2^2+xi3^2)  -- ступінь поляризації
%    велика піввісь ~ cos(chi),  мала ~ |sin(chi)|
%    напрямок обертання = sign(xi2)
% ================================================================

\newcommand{\pR}{0.30}   % радіус нормування вставок (cm)

% ================================================================
%  \polplotStokes{xi1}{xi2}{xi3}{color}
%
%  Малює осі X1-X2 і правильний еліпс поляризації
%  для точки (xi1, xi2, xi3) на сфері Пуанкаре.
% ================================================================
\newcommand{\polplotStokes}[4]{%
  \begin{scope}
    \pgfmathsetmacro{\pp}{sqrt(#1*#1 + #2*#2 + #3*#3)}
    \pgfmathsetmacro{\ppSafe}{max(\pp, 0.001)}
    \pgfmathsetmacro{\psiDeg}{0.5*atan2(#1,#3)}
    \pgfmathsetmacro{\chiDeg}{0.5*asin(#2/\ppSafe)}
    \pgfmathsetmacro{\sma}{\pR*\ppSafe*cos(\chiDeg)}
    \pgfmathsetmacro{\smb}{\pR*\ppSafe*abs(sin(\chiDeg))}
    % осі
    \draw[-{Stealth[length=2.2pt,bend]},thin,black!50]
        ({-1.6*\pR},0)--({1.6*\pR},0)
        node[right,font=\tiny,inner sep=0.5pt]{$X_1$};
    \draw[-{Stealth[length=2.2pt,bend]},thin,black!50]
        (0,{-1.6*\pR})--(0,{1.6*\pR})
        node[above,font=\tiny,inner sep=0.5pt]{$X_2$};
    % еліпс, повернутий на psi
    \begin{scope}[rotate=\psiDeg]
      \pgfmathsetmacro{\isLine}{(\smb < 0.013) ? 1 : 0}
      \ifnum\isLine=1
        \draw[#4, line width=1.2pt](-\sma,0)--(\sma,0);
      \else
        \draw[#4, thick](0,0)ellipse(\sma cm and \smb cm);
        \pgfmathsetmacro{\sensesign}{(#2 >= 0) ? 1 : -1}
        \draw[-{Stealth[length=3.5pt,bend]},#4,thick]
            (\sma,0) arc[start angle=0,
                         end angle={\sensesign*65},
                         x radius=\sma, y radius=\smb];
      \fi
    \end{scope}
  \end{scope}
}

% ================================================================
\tikzset{
  spyglass/.style={
    circle, draw=blue!50, line width=0.7pt,
    inner sep=0pt, fill=white,
    drop shadow={opacity=0.25,
                 shadow xshift=0pt, shadow yshift=-0pt}
  }
}

% ================================================================
%  \addpoint{nodename}{xi1}{xi2}{xi3}{direction}{distance}{color}
%  (викликати всередині scope tdplot_screen_coords)
% ================================================================
\newcommand{\addpoint}[7]{%
  \node[] (#1-lens) at ($(#1)+(#5:#6)$) {%
    \tikz[baseline=0]{\polplotStokes{#2}{#3}{#4}{#7}}%
  };
  \draw[dashed, thin, gray!70, -{Stealth[length=10pt]}] (#1-lens) -- (#1);
}

% ================================================================
%  \addpointangles{nodename}{psi}{chi}{direction}{distance}{color}
%
%  Задаєш точку через фізичні кути еліпса поляризації:
%
%    psi  (-90..90 градусів) -- кут нахилу великої осі
%    chi  (-45..45 градусів) -- кут еліптичності
%                               0   -> лінійна поляризація
%                               ±45 -> кругова поляризація
%
%  Координати точки на сфері (автоматично):
%    xi1 = cos(2*chi) * cos(2*psi)
%    xi2 = sin(2*chi)
%    xi3 = cos(2*chi) * sin(2*psi)
% ================================================================
\newcommand{\addpointangles}[6]{%
  % #1=nodename  #2=psi  #3=chi  #4=direction  #5=distance  #6=color
  \pgfmathsetmacro{\xxi} {cos(2*#3)*cos(2*#2)}
  \pgfmathsetmacro{\xxii} {sin(2*#3)}
  \pgfmathsetmacro{\xxiii}{cos(2*#3)*sin(2*#2)}
  \node[fill=#6!80!black, circle, inner sep=1.5pt]
      (#1) at (\xxi,\xxii,\xxiii){};
  \begin{scope}[tdplot_screen_coords, x=1cm, y=1cm]
    \addpoint{#1}{\xxi}{\xxii}{\xxiii}{#4}{#5}{#6}
  \end{scope}
}

% ================================================================


\tdplotsetmaincoords{68}{115}

\begin{tikzpicture}[tdplot_main_coords, scale=3.6, >=Stealth,
                    every node/.style={font=\small}]

  % ---------------------------------------------------------- СФЕРА
  \shadedraw[tdplot_screen_coords, ball color=blue!12!white, opacity=0.60]
      (0,0) circle (1);

  \tdplotdrawarc[dashed,black,thin]{(0,0,0)}{1}{0}{360}{}{}
  \tdplotdrawarc[black,thin]{(0,0,0)}{1}{-60}{110}{}{}

  \pgfmathsetmacro{\TWOPHI}{0}
  \tdplotsetthetaplanecoords{\TWOPHI}
  \tdplotdrawarc[dashed,black,thin,tdplot_rotated_coords]{(0,0,0)}{1}{0}{360}{}{}
  \tdplotdrawarc[black,thin,tdplot_rotated_coords]{(0,0,0)}{1}{-25}{160}{}{}

  \pgfmathsetmacro{\TWOPHI}{90}
  \tdplotsetthetaplanecoords{\TWOPHI}
  \tdplotdrawarc[dashed,black,thin,tdplot_rotated_coords]{(0,0,0)}{1}{0}{360}{}{}
  \tdplotdrawarc[black,thin,tdplot_rotated_coords]{(0,0,0)}{1}{-35}{140}{}{}

  % ---------------------------------------------------------- ОСІ
  \draw[thick,->](-1.2,0,0)--(2.5,0,0) node[anchor=north east]{$\xi_1$};
  \draw[thick,->](0,-1.1,0)--(0,1.5,0) node[anchor=north west]{$\xi_2$};
  \draw[thick,->](0,0,-1.1)--(0,0,1.5) node[anchor=south]     {$\xi_3$};


  % ==========================================================
  %  ДОВІЛЬНА ТОЧКА
  %  Міняй тільки \myPsi і \myChi -- решта автоматично.
  %
  %  \myPsi (-90..90):  кут нахилу великої осі еліпса
  %  \myChi (-45..45):  кут еліптичності
  %                     0   -> лінійна
  %                     ±45 -> кругова
  % ==========================================================

\begin{pgfonlayer}{background}
\def\myPsi{0}
\foreach \myChi in {0, 11.25, 22.5, ..., {168.75}}
{
 \addpointangles{Sp1}{0}{\myChi}{90}{1.5}{red}
}

%\def\myChi{0}

%\foreach \myPsi in {0, 22.5, 45, ..., 135}
%{
% \addpointangles{Smy}{\myPsi}{\myChi}{0}{1.5}{orange}
%}

\end{pgfonlayer}

%\addpointangles{Sp2}{0}{45}{0}{0.75}{red}
%
%\addpointangles{Sm1}{0}{90}{0}{0.75}{red}
%
%\addpointangles{Sm2}{0}{135}{180}{0.75}{red}

% ----------
%\addpointangles{Sp3}{45}{0}{135}{0.75}{red}
%
%\addpointangles{Sm3}{45}{90}{-45}{0.75}{red}

%\def\myChi{0}
%%\def\myPsi{0}
%\foreach \myPsi in {45, 135}
%{
% \addpointangles{Smy}{\myPsi}{\myChi}{30}{0.5}{orange}
%}
\end{tikzpicture}
\caption{Сфера Пуанкаре: геометричне представлення параметрів Стокса. Координати \((\xi_1, \xi_2, \xi_3)\) визначають точку на сфері для повністю поляризованого світла (\(p=1\)). На рисунку показано як змінюється характер поляризації вздовж екватора.}
\label{pic:Stoke}
\end{figure}
%---------------------------------------------------------
